\documentclass[12pt]{article}
\usepackage[T1]{fontenc}
\usepackage[utf8]{inputenc}
\usepackage{lmodern}
\usepackage{textcomp}
\usepackage{enumitem}
\usepackage{geometry}
\geometry{margin=1in}
\begin{document}
\begin{center}
Answer Key - Biology - Undergraduate Level Assessment 4 \
\small (Answers are ordered exactly as in the question paper.)
\end{center}

\section*{Multiple Choice}
\begin{enumerate}[label=\arabic*.]
\item \textbf{Answer:} A
\\ \textit{Options:} A) DNA migrates from positive charge to negative charge.; B) Smaller DNA travels faster.; C) The DNA migrates only when the current is running.; D) The longer the current is running, the farther the DNA will travel.
\\ \textit{Note:} The correct answer is based on the fundamental principle that DNA, which carries a negative charge due to its phosphate backbone, migrates towards the positive electrode during gel electrophoresis, making the statement about DNA migrating from positive to negative charge incorrect.
\item \textbf{Answer:} A
\\ \textit{Options:} A) hydrogen; B) methane; C) nitrogen; D) sulfur dioxide; E) neon
\\ \textit{Note:} Hydrogen was likely not abundant in the primitive Earth's atmosphere because it is a light gas that would have escaped into space, while other gases such as methane, ammonia, and water vapor were more prevalent in the conditions conducive to the origin of life.
\item \textbf{Answer:} A
\\ \textit{Options:} A) Complete absence of segmentation; B) Hindrance in the development of the nervous system; C) Inability to undergo metamorphosis; D) Absence of a group of contiguous segments; E) Tumor formation in imaginal discs
\\ \textit{Note:} The correct answer is based on the fact that homeotic cluster genes are crucial for establishing the body plan and segment identity during Drosophila development, and their mutation can lead to a failure in proper segmentation, resulting in a complete absence of segmentation.
\item \textbf{Answer:} B
\\ \textit{Options:} A) cross section of muscle tissue; B) longitudinal section of a shoot tip; C) longitudinal section of a leaf vein; D) cross section of a fruit; E) cross section of a leaf
\\ \textit{Note:} A longitudinal section of a shoot tip would most likely provide examples of mitotic cell divisions because the shoot tip is a region of active growth where cells are rapidly dividing to support the development of new plant tissues.
\item \textbf{Answer:} B
\\ \textit{Options:} A) Vinegar has a lower concentration of hydrogen ions than hydroxide ions; B) Vinegar has an excess of hydrogen ions with respect to hydroxide ions; C) Vinegar has a higher concentration of hydroxide ions than water; D) Vinegar has no hydroxide ions; E) Vinegar does not contain hydrogen ions
\\ \textit{Note:} Vinegar is classified as an acid because it has a higher concentration of hydrogen ions (H⁺) compared to hydroxide ions (OH⁻), which results in a lower pH.
\end{enumerate}

\section*{True / False}
\begin{enumerate}[label=\arabic*.]
\setcounter{enumi}{5}
\item \textbf{Answer:} True
\\ \textit{Options:} A) True; B) False
\\ \textit{Note:} The statement is true because increased H+ concentration from lower blood pH promotes the release of O$_{2}$ from hemoglobin, a phenomenon known as the Bohr effect, thereby decreasing hemoglobin's affinity for oxygen.
\item \textbf{Answer:} False
\\ \textit{Options:} A) True; B) False
\\ \textit{Note:} The statement is false because the stomach secretes mucus and bicarbonate to protect its lining from the acidic environment created by gastric acid, rather than high pH levels preventing damage.
\item \textbf{Answer:} False
\\ \textit{Options:} A) True; B) False
\\ \textit{Note:} The flight characteristics of a bird are influenced by both the size and shape of its wings, as wing shape affects aerodynamics, lift, and maneuverability in addition to size.
\item \textbf{Answer:} True
\\ \textit{Options:} A) True; B) False
\\ \textit{Note:} Spindle fibers are indeed composed of microtubules, and they play a crucial role in attaching to chromosomes and ensuring their accurate separation during the process of mitosis.
\item \textbf{Answer:} True
\\ \textit{Options:} A) True; B) False
\\ \textit{Note:} Silk worms exhibit a behavior known as lunar rhythm, where their cocoon spinning is influenced by the phases of the moon, indicating that internal biological processes can synchronize with external environmental cues.
\end{enumerate}

\section*{Long Form Response}
\begin{enumerate}[label=\arabic*.]
\setcounter{enumi}{10}
\item \textbf{Answer:} Stabilizing selection is an evolutionary process that favors average traits within a population, leading to a decrease in variation for those traits. In the context of human birth weight, stabilizing selection operates by favoring infants that are around the average weight of 3 kg. Infants of this weight are more likely to survive and thrive, as they are less prone to complications associated with being either underweight (which can lead to higher mortality rates) or overweight (which can result in delivery complications). This selection pressure enhances the evolutionary fitness of the average weight, as it increases the likelihood of successful reproduction and survival of offspring, thereby maintaining the population's overall stability in birth weight.
\\ \textit{Note:} correct because it accurately describes how stabilizing selection promotes the survival of infants with an average birth weight of 3 kg by reducing the likelihood of negative health outcomes associated with both low and high birth weights, thereby enhancing evolutionary fitness.
\item \textbf{Answer:} In biology, a "society" can be exemplified by a flock of birds, which represents a basic form of social organization known as a motion group. Within this group, birds engage in various forms of communication, such as vocalizations and body movements, to coordinate their movements and maintain cohesion while traveling. This cooperative behavior is crucial for their survival, as it enhances their ability to detect predators and respond effectively to threats. By working together, the flock can employ collective vigilance, allowing individual birds to focus on foraging while still benefiting from the safety in numbers. Overall, the social interactions within a flock illustrate how communication and cooperation are vital for enhancing survival in the wild.
\\ \textit{Note:} The answer correctly identifies a flock of birds as a biological society, emphasizing the importance of their communication and cooperative behaviors in enhancing their survival by improving predator detection and response.
\end{enumerate}

\vfill
\noindent\textit{End of Answer Key}
\end{document}